\chapter*{Conclusion générale et perspectives}
\addcontentsline{toc}{chapter}{Conclusion générale et perspectives}
\markboth{Conclusion générale}{Conclusion générale}

Ce stage, réalisé au sein de \textbf{\entreprise}, avait pour objectif de concevoir et réaliser
une plateforme de supervision et d'alerte auto-hébergée, simple à installer et fiable. Le point
de départ était un constat vécu : mettre en place Prometheus, Alertmanager et Grafana demande
de configurer et de synchroniser plusieurs outils, ce qui est long et fragile pour une petite
équipe.

\textbf{SentinelOps} répond à ce besoin. Une seule commande, \code{sen run}, construit et démarre
l'ensemble et affiche \og Ready \fg{} lorsque tout fonctionne. L'outil collecte les métriques de
l'hôte, des conteneurs et de cibles Prometheus, détecte les incidents par des règles avec durée
et par un score z robuste, regroupe les alertes liées en incidents, explique chaque alerte par
des preuves mesurées et un instantané des processus, et notifie les équipes sur Slack ou par
courriel sans jamais perdre une notification.

Au-delà des fonctionnalités, le travail a surtout porté sur la qualité. Une revue du prototype
initial a relevé seize problèmes, de la base de données définie à trois endroits à une API
entièrement ouverte. La version 0.2 les corrige par des choix d'architecture : schéma unique par
migrations, automate d'états, outbox transactionnelle, API fermée par défaut, accès à Docker au
moindre privilège et conteneurs durcis. Cinquante et un tests automatisés et une chaîne
d'intégration continue avec analyse de vulnérabilités vérifient chaque modification.

Sur le plan personnel, ce stage m'a permis de \afaire{compléter : compétences acquises,
par exemple la conception d'un système concurrent autour d'une base partagée, les pratiques
DevSecOps, la rigueur des tests, le travail en autonomie avec un encadrant}.

\section*{Perspectives}

Plusieurs évolutions sont prévues dans la feuille de route du projet :
\begin{itemize}
  \item \textbf{détection par apprentissage automatique} : ajouter un détecteur Isolation
        Forest multivarié, sans modifier le gestionnaire d'alertes grâce à l'interface
        \code{Detector} ;
  \item \textbf{indicateurs de fiabilité} : calculer le MTTD et le MTTR à partir des dates déjà
        enregistrées pour chaque alerte ;
  \item \textbf{supervision de plusieurs hôtes} : des agents distants qui envoient leurs
        métriques au dashboard ;
  \item \textbf{PostgreSQL} en option pour les volumes plus importants ;
  \item \textbf{déploiement automatisé} avec Terraform et un chart Helm pour Kubernetes ;
  \item \textbf{assistant d'analyse} : un modèle de langage exécuté localement (par exemple via
        Ollama) pour résumer un incident et proposer des pistes de résolution ;
  \item \textbf{HTTPS intégré et rôles} (lecture seule, opérateur, administrateur).
\end{itemize}
